% page 122 of the facsimile (image p-121 of the scan)
% block 167.6 x 254.9 mm ; left margin 8.0 mm
\pgc{-12.700mm}{0.000mm}{
\l{\cel{7}ll2.}
\l{}
\l{}
\l{\sou{diplasar}. (\sou{trans.}) Deprenar (ulo, ulu) de lua plaso. - EF.}
\l{}
\l{\sou{diploo}.\cel{2}(\sou{anat}.) I.Asemblajo de du lami osta, kompakta, ye}
\l{\cel{3}qui konsistas la surfaci interna ed extera di la kranio.}
\l{\cel{3}- II. Tisuo sponjatra quan onu renkontras inter la surfaci}
\l{\cel{3}di ta osti, e, generale, en la osti plata.}
\l{}
\l{\sou{diplomo}. I. (\sou{Olim)}. Akto oficala (faldita e siglita) qua emanis}
\l{\cel{3}de suvereno ed establisis o konfirmis yuro, privilejo. - II.}
\l{\cel{3}(\sou{Nun}). Akto emananta de universitato, fakultato, skolo publi-}
\l{\cel{3}ka, qua grantas grado, titulo. - DEFIRS.}
\l{}
\l{\sou{diplomaco}. Fako di politiko, qua koncernas la relati di populo,}
\l{\cel{3}di guvernistaro, kun altra, la negocii, la kontrati interna-}
\l{\cel{3}ciona. - DEFIRS.}
\l{}
\l{\sou{diplomatiko}.\cel{2}Cienco di qua la studiajo esas la diplomi, charti,}
\l{\cel{3}ed altra tituli, la verifiko di la autentikeso, di la integreso}
\l{\cel{3}di la akti, e c. - DEFIRS.}
\l{}
\l{\sou{diplostemono}. (\sou{bot.}) Floro o planto di qua la quanto di la sta-}
\l{\cel{3}mini esas duopla ye olta di la sepali e di la petali.\cel{2}-}
\l{}
\l{\sou{diporto}. (\sou{Financo).}\cel{2}Ajorno, til liquidaco borsala proxima, di}
\l{\cel{3}la livro di valor-akti (acioni, obligacioni), vendita, pag-}
\l{\cel{3}ante premio ; la premio pagata por ta ajorno. - DFIS.}
\l{}
\l{\sou{diptero}.\cel{2}(\sou{zool.)}\cel{2}Insekto sugera, di qua nur la du ali supra}
\l{\cel{3}esas developita, le infra esante remplasigita da organi rudi-}
\l{\cel{3}menta. - DEFIS.}
\l{}
\l{\sou{diptiko}.\cel{2}I. (En\cel{1}\sou{la epoki antiqua})\cel{2}Asemblajo de du tabuleti}
\l{\cel{3}ivora sur qua onu skribis (en Roma) la nomo di la konsuli e}
\l{\cel{3}di la stat-oficieri alta-ranga ; en la Eklezio katolika, la}
\l{\cel{3}nomi di la episkopi, di la bonfaceri di la Eklezio, di la}
\l{\cel{3}martiri, e c. - II. (\sou{Nun}). Pikturo basreliefa, kovrita da}
\l{\cel{3}quaza shutro di qua la surfaco interna esas anke pikturoza}
\l{\cel{3}e skulturoza. - DEFIS.}
\l{}
\l{\sou{direciono}. Lineo rekta segun qua korpo su movas o pozesas.-DEFIRS.}
\l{}
\l{\sou{direktar}. (\sou{trans.})\cel{2}I. Funcionigar segun ta o ca konduto-lineo ;}
\l{\cel{3}(aludante kleriko) trasar ad ulu la konduto-lineo quan lu de-}
\l{\cel{3}vas sequar pri lua koncienco. - II. Movigar ad skopo indikita.}
\l{\cel{3}- DEFIRS.}
\l{}
\l{\sou{Direktorio.}\cel{2}(\sou{Historio di Francia}). Komisitaro de kin membri,}
\l{\cel{3}qua investesis per la povo exekutala, de l795 til l799).-}
\l{\cel{3}DEFIRS.}
\l{}
\l{direta.(\sou{pri}\cel{1}treno). Qua, segun lua itinerario, atingas suades-}
\l{\cel{3}tino-staciono sen haltir ye la multa stacioni di la inter-}
\l{\cel{3}spaco. II. (\sou{pri voyo)}. Rekta ; sen-sinua o poke-sinua. III.}
\l{\cel{3}(\sou{pri cito}). En qua onu raportas, uzante la pronomo \textquotedbl{}me\textquotedbl{}, la}
\l{\cel{3}vorti pronuncita (reale o konjektate) da ulu. IV. (\sou{pri metodo})}
\l{\cel{3}Qua iras rekte a la skopo. - DEFIS.}
}
